\documentclass[12pt]{article}
\usepackage{fullpage}
\usepackage{amsmath,amssymb,amsthm}

\newtheorem{theorem}{Theorem}
\newtheorem{lemma}{Lemma}

\newcommand{\OP}{\operatorname{OP}}
\newcommand{\inv}{\operatorname{inv}}
\newcommand{\Hilb}{\operatorname{Hilb}}
\newcommand{\Stir}{\operatorname{Stir}}

\begin{document}

\begin{center}
{\bf Hook coefficients in the universal Hilbert series of diagonal coinvariants}
\end{center}

\section*{Problem statement and interpretation}

I interpret \(R_n^{(m)}\) as the diagonal coinvariant algebra
\[
 R_n^{(m)}=\mathbb C[x_i^{(r)}:1\le i\le n,\ 1\le r\le m]\,
/\,\langle \mathbb C[x_i^{(r)}]^{\mathfrak S_n}_{+}\rangle,
\]
where \(\mathfrak S_n\) acts by simultaneously permuting the subscript \(i\) in all
\(m\) sets of variables.  Thus the coefficients \(c_\lambda(n)\) are the stable
multiplicities of the polynomial \(GL_m\)-irreducibles \(S_\lambda(\mathbb C^m)\).

\section*{The combinatorial model}

For \(r,k\ge 1\), let \(\OP_{r,k}\) be the set of ordered set partitions
\(\sigma=(B_1|\cdots|B_k)\) of \([r]\) into \(k\) nonempty blocks.  An inversion of
\(\sigma\) is a pair \((i,j)\) such that
\[
        i<j,\qquad i=\min B_m,\qquad j\in B_\ell,\qquad \ell<m .
\]
Let \(\inv(\sigma)\) be the number of such pairs.  When \(k=r\), this is the
usual inversion number of the permutation obtained by reading the singleton
blocks from left to right.

\begin{theorem}
Let \(\lambda=(a,1^b)\) be a hook, with \(a\ge 1\) and \(b\ge 0\).  Then
\(c_{(a,1^b)}(n)\) is the number of ordered set partitions
\(\sigma=(B_1|\cdots|B_k)\) of \([r]\), for some \(r\le n\), such that
\[
        r-k=b+1,\qquad r\hbox{ is not a singleton block of }\sigma,\qquad
        \inv(\sigma)=a-1 .
\]
Equivalently,
\[
\sum_{a\ge1,\ b\ge0} c_{(a,1^b)}(n)\,q^{a-1}z^b
=
\sum_{\substack{2\le r\le n\\1\le k\le r-1}}
\ \sum_{\substack{\sigma\in \OP_{r,k}\\
                  r\text{ not singleton in }\sigma}}
        q^{\inv(\sigma)}z^{r-k-1}.
\]
The empty partition coefficient is \(c_{\varnothing}(n)=1\).
\end{theorem}

\begin{proof}
We use two standard inputs.  First, diagonal supersymmetry says that the same
universal coefficients \(c_\lambda(n)\) occur if one replaces all but one
bosonic set of variables by one fermionic set.  Thus the Hilbert series of the
superspace coinvariant ring
\[
 SR_n=\bigl(\mathbb C[x_1,\ldots,x_n]\otimes
          \bigwedge(\theta_1,\ldots,\theta_n)\bigr)
/\langle(\mathbb C[x]\otimes\bigwedge\theta)^{\mathfrak S_n}_{+}\rangle
\]
has the expansion
\[
        \Hilb(SR_n;q,z)=\sum_\lambda c_\lambda(n)s_\lambda(q/z),
\]
where \(q\) records \(x\)-degree and \(z\) records \(\theta\)-degree.  In one
bosonic and one fermionic variable, \(s_\lambda(q/z)\) vanishes unless
\(\lambda\) is a hook, and
\[
        s_{(a,1^b)}(q/z)=q^a z^b+q^{a-1}z^{b+1}
        =q^{a-1}z^b(q+z).
\]
Therefore
\[
        \Hilb(SR_n;q,z)
        =1+(q+z)\sum_{a\ge1,\ b\ge0}
              c_{(a,1^b)}(n)q^{a-1}z^b .          \tag{1}
\]

Second, the Hilbert series of \(SR_n\) is
\[
        H_n(q,z):=\Hilb(SR_n;q,z)
        =\sum_{k=1}^{n} z^{\,n-k}[k]!_q\,\Stir_q(n,k),                 \tag{2}
\]
where \([k]_q=1+q+\cdots+q^{k-1}\), \([k]!_q=[k]_q[k-1]_q\cdots[1]_q\), and
\[
        \Stir_q(n,k)=\Stir_q(n-1,k-1)+[k]_q\Stir_q(n-1,k).
\]
Put \(F_{r,k}=[k]!_q\Stir_q(r,k)\).  From the recurrence,
\[
        F_{r,k}=[k]_qF_{r-1,k-1}+[k]_qF_{r-1,k}.
\]
Using \([k+1]_q=1+q[k]_q\) in (2), we get
\[
\begin{aligned}
H_n(q,z)
&=\sum_{k=1}^{n}z^{n-k}\bigl([k]_qF_{n-1,k-1}
                         +[k]_qF_{n-1,k}\bigr)\\
&=H_{n-1}(q,z)
 +(q+z)\sum_{k=1}^{n-1}[k]_qF_{n-1,k}z^{n-1-k}.
\end{aligned}
\]
Since \(H_1(q,z)=1\), iteration gives
\[
        \frac{H_n(q,z)-1}{q+z}
        =
        \sum_{r=1}^{n-1}\sum_{k=1}^{r}
             [k]_q\,F_{r,k}\,z^{r-k}.                    \tag{3}
\]

It remains to identify the right side of (3).  The polynomial \(F_{r,k}\) is
the inversion enumerator of ordered set partitions:
\[
        F_{r,k}=\sum_{\tau\in \OP_{r,k}}q^{\inv(\tau)}.                \tag{4}
\]
Indeed, both sides satisfy the same recurrence.  To see this directly, insert
the largest letter \(r\) into an ordered partition of \([r-1]\).  If \(r\) is
inserted as a singleton block, or into an existing block, at position \(i\) from
the left, the new inversions involving \(r\) are exactly the minima of the
blocks to its right, so their number is \(k-i\).  Summing over \(i=1,\ldots,k\)
contributes \([k]_q\), giving the recurrence above and the initial condition.

Now restrict this insertion argument to insertions of \(r+1\) into an existing
block of an ordered partition \(\tau\in\OP_{r,k}\).  The result is precisely an
ordered partition \(\sigma\in\OP_{r+1,k}\) in which \(r+1\) is not a singleton,
and the extra inversion contribution is again one of \(0,1,\ldots,k-1\).
Hence
\[
 [k]_qF_{r,k}
 =
 \sum_{\substack{\sigma\in\OP_{r+1,k}\\ r+1\text{ not singleton}}}
      q^{\inv(\sigma)} .
\]
Substituting this into (3), and comparing with (1), gives exactly
\[
\sum_{a,b}c_{(a,1^b)}(n)q^{a-1}z^b
=
\sum_{\substack{2\le r\le n\\1\le k\le r-1}}
\sum_{\substack{\sigma\in \OP_{r,k}\\r\text{ not singleton}}}
q^{\inv(\sigma)}z^{r-k-1}.
\]
Extracting the coefficient of \(q^{a-1}z^b\) proves the claimed interpretation.
\end{proof}

\begin{thebibliography}{9}

\bibitem{Bergeron}
F. Bergeron,
\emph{Multivariate diagonal coinvariant spaces for complex reflection groups},
Adv. Math. 239 (2013), 97--108.

\bibitem{Lentfer}
J. Lentfer,
\emph{Diagonal supersymmetry for coinvariant rings},
arXiv:2505.14885 (2025).

\bibitem{RhoadesWilson}
B. Rhoades and A. T. Wilson,
\emph{The Hilbert series of the superspace coinvariant ring},
Forum Math. Pi 12 (2024), e16.

\end{thebibliography}

\end{document}
